\section{Conclusion}\label{sec:conclusion}
This study built two simulators of the same room for Pepper, one from a photogrammetric scan and
one by hand from tape measurements, trained a SAC navigation policy in each under identical domain
randomisation, and deployed both on the real robot without further training.

On H1, the photogrammetry policy has a significantly smaller sim-to-real gap: its success rate fell
by 8 percentage points from simulation to the real robot, against 37 for
the standard policy. On H2, its odometry-estimated success is 9 of 20 real routes against 6, which
points the same way but is not statistically significant. Most real failures of both policies
happened at the couch and the bag, where collision protection, unseen in training, stopped the
robot.

For the two policies tested, the results suggest that a reconstructed room transfers to the real
robot more faithfully than a hand-built one made from the same measurements. Whether this holds
across training seeds needs more than one deployed policy per branch. Photogrammetry alone is not enough, however:
transfer also depends on modelling the robot's own behaviour, such as its collision protection,
and confirming the real-world difference needs larger trials with external position tracking.
